\documentclass{article}
\usepackage{../math_template}

\author{Jonathan Kasongo}
\title{Australian Math Olympiad 2020 --- P2/4}

\begin{document}
\maketitle

\begin{problem}{Australian Math Olympiad 2020 --- P2/4}
Amy and Bec play the following game. Initially, there are three piles, each containing
2020 stones. The players take turns to make a move, with Amy going first. Each
move consists of choosing one of the piles available, removing the unchosen pile(s)
from the game, and then dividing the chosen pile into 2 or 3 non-empty piles. A
player loses the game if they are unable to make a move.
\vspace{0.2cm}

Prove that Bec can always win the game, no matter how Amy plays.
\end{problem}

\begin{solution}{Solution}
The following claim is the heart of the problem. \\

\textit{Claim 1: If, on Bec's turn, there is a pile of
$x$ stones, where} $x \nequiv 1 \ \text{(mod } 3)$
\textit{then, in the subsequent turns she is able to force Amy to
always choose a pile of $y < x$ stones, where} $y \equiv 1 \ \text{(mod } 3
)$. \\

\textit{Proof:} We have two cases,

\begin{itemize}
\item If $x \equiv 0 \ \text{(mod } 3)$, then Bec can split the pile of
$x$ stones like so: $(1, 1, x-2)$. Each of these numbers is congruent to 1
mod 3.
\item If $x \equiv 2  \ \text{(mod } 3)$, then Bec can split the pile of
$x$ stones like so: $(1, x-1)$. Each of these numbers is congruent to 1 mod
3.
\end{itemize}

Thus on Amy's next turn she must pick a pile with $y$ stones, with
$y \equiv 1 \ \text{(mod } 3)$. Also since Amy can never choose a pile of
stones with 1 stone, it must be the case that $y \in \{ x-2, x-1 \}$
and so $y < x$.
Now when Amy makes her move,
she will split the pile such that there is at least one new pile with
$z \nequiv 1 \ \text{(mod } 3)$ stones. Suppose, FTSOC, that there isn't
such a pile.

\begin{itemize}
\item If Amy splits the pile of $y$ stones into 2 piles with $a,b$ stones
such that $a,b \equiv 1 \ \text{(mod } 3)$, then $2 \equiv a+b = y
\equiv 1 \ \text{(mod } 3)$, a contradiction.
\item If Amy splits the pile of $y$ stones into 3 piles with $a,b,c$ stones
such that $a,b,c \equiv 1 \ \text{(mod } 3)$, then $0 \equiv a+b+c = y
\equiv 1 \ \text{(mod } 3)$, a contradiction.
\end{itemize}

Thus after this point, Bec can always choose a pile of $x$ stones where
$x \nequiv 1 \ \text{(mod } 3)$ and in the subsequent turns she is also
able to force Amy to always choose a pile of $y$ stones,
where $y \equiv 1 \ \text{(mod } 3)$.
$\square$ \\

Since $2020 = 673 \cdot 3 + 1 \equiv 1 \ \text{(mod } 3)$. Due to claim 1,
Bec has a strategy to force Amy to always choose piles of stones with
$k \equiv 1 \ \text{(mod } 3)$, whilst she can always avoid choosing piles
of stones with $m \nequiv 1 \ \text{(mod } 3)$.
Again due to claim 1, that means that Amy will eventually
have to pick a pile of stones with 1 stone in that pile, but since it is
impossible to divide this pile into 2 or 3 non-empty piles, Bec will win.
$\blacksquare$
\end{solution}



\end{document}
